% ═══════════════════════════════════════════════════════════════════════════════
%  CHAPITRE 3 : ANALYSE DES BESOINS ET CONCEPTION DU SYSTÈME
% ═══════════════════════════════════════════════════════════════════════════════
\newpage
\section{Analyse des besoins et conception du système}
\label{ch3}

\subsection{Introduction}
\label{ch3:intro}

Ce chapitre détaille la phase de conception du système de gestion des soutenances. Il définit les exigences fonctionnelles et non fonctionnelles, formalise les règles de gestion qui régissent la planification, et présente la modélisation UML du système. Cette étape est cruciale pour garantir que la solution technique répond précisément aux besoins du département. Dans un projet de développement logiciel, la qualité de la phase de conception détermine directement la maintenabilité et l'évolutivité du système résultant. C'est pourquoi nous avons adopté une démarche rigoureuse combinant analyse des besoins, modélisation UML et définition formelle des règles de gestion, afin de disposer d'un référentiel clair avant de passer à l'implémentation.

\subsection{Présentation générale du projet}
\label{ch3:presentation_generale}

Le projet consiste en la réalisation d'une application web centralisée permettant l'organisation complète des sessions de soutenances. Contrairement à une simple gestion d'agenda, le système agit comme un moteur de planification intelligent capable d'orchestrer des sessions globales (normale, rattrapage) avec des paramètres configurables incluant la durée des soutenances, les pauses entre les passages, et les coefficients d'évaluation par rôle.

Le système couvre l'ensemble du cycle de vie d'une soutenance : de la création des projets et la constitution des jurys, en passant par la planification interactive des créneaux via un éditeur visuel par glisser-déposer (\textit{Defense Designer}), jusqu'à la saisie des évaluations et la génération automatique des documents administratifs (convocations, procès-verbaux, fiches d'évaluation). Chaque étape est validée par un moteur de détection de conflits garantissant l'intégrité du planning.

Le périmètre du projet a été défini en collaboration avec l'encadrant et le département d'accueil. Sont inclus dans le périmètre : la gestion des utilisateurs et des rôles, la configuration académique (filières, niveaux, salles), la gestion des projets et des groupes d'étudiants, la constitution des jurys, la planification interactive des créneaux, la détection automatique de conflits, la saisie des évaluations, la génération de documents et le système de notifications. Sont exclus du périmètre : la gestion financière, la publication des résultats finaux (bulletins), l'intégration avec le système d'information universitaire existant, et le développement d'une application mobile.

\subsection{Cahier des charges}
\label{ch3:cahier_charges}

\subsubsection{Acteurs du système}
\label{ch3:acteurs}

Le système repose sur une gestion stricte des accès basée sur les rôles (\textit{Role-Based Access Control — RBAC}). Quatre acteurs principaux sont identifiés, chacun disposant d'un périmètre fonctionnel spécifique. Le choix de quatre rôles uniquement, sans possibilité de superposition, simplifie considérablement la logique de contrôle d'accès tout en couvrant l'ensemble des besoins identifiés.

L'\textbf{Administrateur} est responsable de la configuration structurelle de l'établissement. Il gère les départements, filières, niveaux et salles. Il est également en charge de la création des comptes utilisateurs (import CSV, activation) et de la configuration globale des paramètres du système (durée des soutenances, templates de jurys, configuration email). Ce rôle est le seul à pouvoir accéder à la configuration système et aux journaux d'audit.

Le \textbf{Coordinateur} est l'acteur central du système. Il gère l'ensemble du cycle de planification : import et validation des projets, constitution des jurys selon les templates de rôles, utilisation du \textit{Defense Designer} pour la planification interactive des créneaux, consultation du tableau de bord des conflits, et suivi des évaluations. Il peut également publier le planning final et générer les documents administratifs. Ce rôle nécessite une vue d'ensemble de toutes les soutenances et une compréhension des contraintes de chaque acteur.

L'\textbf{Enseignant} est membre du corps professoral impliqué dans les soutenances. Il consulte son planning personnel, déclare ses indisponibilités via un calendrier interactif, saisit les notes des étudiants qui lui sont assignés, et accède à ses convocations. Son interaction avec le système est principalement consultative, à l'exception de la saisie des notes et de la déclaration d'indisponibilités.

L'\textbf{Étudiant} est l'utilisateur final du système. Il visualise les informations relatives à sa soutenance (date, salle, jury), consulte la composition de son groupe de projet, et télécharge sa convocation officielle au format A4. Son périmètre d'action est le plus limité, ce qui est cohérent avec son rôle dans le processus de soutenance.

Le tableau~\ref{tab:acteurs_use_cases} résume la répartition des fonctionnalités principales par acteur.

\begin{table}[H]
    \centering
    \scriptsize
    \begin{tabular}{lp{2.4cm}p{2.4cm}p{2.4cm}p{2.4cm}}
        \toprule
        \rowcolor{primary!10}
        \textbf{Fonctionnalité} & \textbf{Admin} & \textbf{Coord.} & \textbf{Enseign.} & \textbf{Étudiant} \\
        \midrule
        Gestion départements & \checkmark & — & — & — \\
        Gestion salles & \checkmark & — & — & — \\
        Gestion utilisateurs & \checkmark & — & — & — \\
        Configuration système & \checkmark & — & — & — \\
        Gestion projets & — & \checkmark & — & — \\
        Constitution jurys & — & \checkmark & — & — \\
        Planification créneaux & — & \checkmark & — & — \\
        Détection conflits & — & \checkmark & — & — \\
        Saisie indisponibilités & — & — & \checkmark & — \\
        Saisie évaluations & — & — & \checkmark & — \\
        Consultation planning & — & \checkmark & \checkmark & \checkmark \\
        Consultation groupe & — & — & — & \checkmark \\
        Téléchargement documents & — & — & — & \checkmark \\
        \bottomrule
    \end{tabular}
    \caption{Répartition des fonctionnalités par acteur}
    \label{tab:acteurs_use_cases}
\end{table}

\subsubsection{Besoins fonctionnels}
\label{ch3:besoins_fonctionnels}

Les besoins sont organisés par modules métier, chacun correspondant à un sous-système cohérent de l'application. Chaque module a été identifié lors des entretiens avec le coordinateur et l'encadrant, puis priorisé en fonction de son impact sur le processus global de soutenance.

Le \textbf{module Gestion Identité} assure l'authentification sécurisée via JWT avec tokens d'accès et de rafraîchissement. Il gère la récupération de mot de passe avec lien à durée limitée (1 heure), l'activation de compte avec validation de la force du mot de passe (via la bibliothèque zxcvbn), et l'interface de connexion unifiée avec redirection automatique basée sur le rôle de l'utilisateur. Ce module constitue la porte d'entrée du système et doit garantir la confidentialité des données tout en offrant une expérience fluide.

Le \textbf{module Gestion Académique} couvre l'administration des filières (MIP, BCG, etc.), des niveaux (L3, M1, M2) et des grades académiques. Il inclut l'import massif des étudiants et des salles via fichiers CSV avec validation des données et gestion des erreurs ligne par ligne. L'interface CRUD complète avec recherche et pagination permet de maintenir les référentiels à jour tout au long de l'année universitaire.

Le \textbf{module Gestion des Projets et Jurys} gère la création et le suivi des projets avec des statuts (proposé, approuvé, rejeté). Il permet la constitution des jurys selon des templates configurables avec rôles normalisés (Président, Rapporteur, Examinateur), et assure le suivi de la charge des enseignants pour éviter les sur-affectations. Ce module est le point de départ du processus de planification, car les projets et jurys sont les entités qui seront ensuite placés dans le planning.

Le \textbf{module Planification (Defense Designer)} est l'élément le plus innovant du système. Il fournit une interface interactive de type canevas permettant au coordinateur de planifier les soutenances par glisser-déposer. Les jurys sont listés dans une barre latérale et peuvent être déposés sur un calendrier dynamique organisé par salles et créneaux horaires. L'état local est maintenu côté client et synchronisé avec le serveur uniquement lors de la sauvegarde explicite. Un algorithme d'auto-génération permet de peupler automatiquement le planning en optimisant l'occupation des salles.

Le \textbf{module Moteur de Conflits} assure la validation systématique de 8 règles de gestion lors de chaque modification du planning (voir \ref{ch3:regles_gestion}). Le moteur est dupliqué côté client pour un retour instantané lors des interactions et exécuté côté serveur pour garantir l'intégrité lors de la sauvegarde. Cette architecture double-validation est un choix architectural majeur qui combine réactivité et fiabilité.

Le \textbf{module Évaluation} gère la saisie des notes par rôle de jury avec des coefficients dynamiques configurés par session. Le calcul automatique de la moyenne pondérée élimine les erreurs de calcul manuel. Le suivi de l'avancement des évaluations avec indication visuelle du statut (en attente, complétée) permet au coordinateur d'identifier rapidement les évaluations en retard.

Le \textbf{module Édition Documentaire} génère 5 types de documents (convocation jury, convocation étudiant, fiche d'évaluation, liste de présence, procès-verbal) via des vues React optimisées pour l'impression A4 utilisant la stratégie \texttt{@media print}. Cette approche client-side évite la complexité d'une génération PDF serveur tout en produisant des documents de qualité professionnelle.

Le \textbf{module Notifications et Audit} fournit un système de notifications internes avec différents niveaux (success, warning, error) et la journalisation des actions critiques avec piste d'audit complète pour la traçabilité. Chaque action significative (création, modification, suppression) est enregistrée avec l'identifiant de l'utilisateur, l'horodatage et les détails de l'opération.

\subsubsection{Besoins non fonctionnels}
\label{ch3:besoins_non_fonctionnels}

Les besoins non fonctionnels définissent les qualités transverses du système, souvent aussi critiques que les fonctionnalités elles-mêmes pour garantir l'adoption et la satisfaction des utilisateurs.

La \textbf{réactivité} est un besoin fondamental pour le Defense Designer. Le retour sur conflit lors du drag-and-drop doit être instantané (latence $< 100$ms) grâce à la validation locale via le moteur de conflits TypeScript. L'interface utilisateur ne doit pas présenter de blocage perceptible lors des interactions, ce qui impose de maintenir l'état du planning dans la mémoire du navigateur et de ne jamais effectuer d'appels réseau pendant les opérations de glisser-déposer.

La \textbf{sécurité} repose sur le chiffrement des mots de passe via BCrypt, la protection des endpoints par tokens JWT avec validation côté serveur (Spring Security), et le contrôle d'accès basé sur les rôles via annotations \texttt{@PreAuthorize}. Les tokens ont une durée de validité limitée (2 heures) et sont stockés dans le \texttt{localStorage} du navigateur. Une attention particulière a été portée à la protection contre les attaques par injection SQL, les failles XSS et les erreurs d'autorisation horizontale.

La \textbf{disponibilité} est assurée par une architecture stateless du backend, permettant un déploiement horizontal scalable. L'API REST ne maintient pas d'état de session côté serveur, chaque requête étant autonome via le token JWT. Cette conception facilite également les tests et le débogage, car chaque requête contient toute l'information nécessaire à son traitement.

L'\textbf{ergonomie} vise une interface moderne et responsive basée sur Tailwind CSS 4 et shadcn/ui. La navigation est contextualisée par rôle avec une sidebar adaptative qui ne présente que les fonctionnalités accessibles à l'utilisateur connecté. Le thème clair/sombre avec persistance du choix améliore le confort d'utilisation lors de sessions de travail prolongées.

La \textbf{maintenabilité} est garantie par une architecture en couches stricte (Contrôleur $\rightarrow$ Service $\rightarrow$ Repository $\rightarrow$ Entité) avec séparation claire des responsabilités. L'utilisation de DTOs pour l'isolation des couches API et de MapStruct pour le mapping objet protège le modèle interne des variations de l'interface exposée.

La \textbf{qualité} est mesurée par une couverture de code minimale de 85\% (lignes) et 70\% (branches) via JaCoCo côté Java. Le linting automatisé avec ESLint (TypeScript) et Spotless/Checkstyle/PMD/SpotBugs (Java) intégré dans le cycle Maven assure la conformité du code aux conventions de l'équipe.

\subsection{Règles de gestion}
\label{ch3:regles_gestion}

Le moteur de conflits implémente les règles de gestion (RG) suivantes, garantissant l'intégrité du planning. Ces règles ont été définies en collaboration avec le coordinateur du département MIP et formalisées sous forme de contraintes vérifiables algorithmiquement.

\begin{projectbox}{Règles de Gestion du Planning}
\begin{description}
    \item[RG1 (Enseignants)] Un enseignant ne peut être affecté à deux soutenances dont les horaires se chevauchent.
    \item[RG2 (Salles)] Une salle ne peut accueillir qu'une seule soutenance pour un créneau donné.
    \item[RG3 (Capacité)] Le nombre d'étudiants d'un projet ne doit pas excéder la capacité de la salle assignée.
    \item[RG4 (Période)] Toutes les soutenances d'une session doivent être planifiées entre la \textit{startDate} et la \textit{endDate} de ladite session.
    \item[RG5 (Encadrants)] Un encadrant ne doit pas être sollicité sur plusieurs projets le même jour pour éviter la surcharge.
    \item[RG6 (Unicité)] Un projet ne peut être planifié qu'une seule fois par session.
    \item[RG7 (Pause)] Un intervalle minimum (\textit{breakDuration}) doit être respecté entre deux soutenances dans une même salle.
    \item[RG8 (Disponibilité)] Aucune soutenance ne peut être planifiée sur un créneau déclaré comme indisponible par un membre du jury.
\end{description}
\end{projectbox}

Ces règles sont implémentées dans deux moteurs distincts mais fonctionnellement identiques : un moteur Java côté serveur (\texttt{ConflictDetectionService}) qui valide lors de la sauvegarde, et un moteur TypeScript côté client (\texttt{conflict-engine.ts}) qui fournit un retour instantané lors des interactions dans le Defense Designer.

Le tableau~\ref{tab:rg_mapping} associe chaque règle à son ou ses implémentations et à la sévérité de la violation. Les règles classées « error » bloquent l'opération, tandis que les règles classées « warning » affichent une alerte mais autorisent l'enregistrement.

\begin{table}[H]
    \centering
    \small
    \begin{tabular}{llccc}
        \toprule
        \rowcolor{primary!10}
        \textbf{Règle} & \textbf{Description} & \textbf{Java} & \textbf{TypeScript} & \textbf{Sévérité} \\
        \midrule
        RG1 & Chevauchement enseignant & \checkmark & \checkmark & error \\
        RG2 & Occupation de salle & \checkmark & \checkmark & error \\
        RG3 & Capacité de salle & \checkmark & \checkmark & error \\
        RG4 & Bornes de période & \checkmark & \checkmark & error \\
        RG5 & Surcharge encadrant & \checkmark & \checkmark & warning \\
        RG6 & Unicité de projet & \checkmark & \checkmark & error \\
        RG7 & Intervalle de pause & \checkmark & \checkmark & error \\
        RG8 & Indisponibilité & \checkmark & \checkmark & error \\
        \bottomrule
    \end{tabular}
    \caption{Correspondance entre les règles de gestion et leurs implémentations}
    \label{tab:rg_mapping}
\end{table}

Lorsqu'une règle ne peut être satisfaite (par exemple, aucun créneau disponible pour un jury donné), le système ne propose pas de mécanisme d'override automatique. Le coordinateur doit soit modifier les paramètres de la session (dates, salles), soit reconfigurer la composition du jury, soit déclarer une indisponibilité compensatoire. Cette approche conservative garantit l'intégrité du planning au prix d'une flexibilité réduite, un choix délibéré compte tenu du caractère critique des contraintes académiques.

\subsection{Conception UML}
\label{ch3:conception_uml}

La modélisation UML permet de formaliser les différentes visions du système : qui fait quoi (cas d'utilisation), comment les données sont structurées (classes), comment les composants interagissent (séquence), et comment les processus se déroulent (activité). Nous avons utilisé les diagrammes UML 2.5, implémentés via PlantUML pour faciliter la maintenance et la régénération automatique des figures.

\subsubsection{Diagrammes de cas d'utilisation}
\label{ch3:uc}

Les diagrammes de cas d'utilisation délimitent les responsabilités de chaque acteur. Ils ont été construits à partir des entretiens avec le coordinateur et l'encadrant, puis validés lors des revues de conception.

Le diagramme de la \ref{fig:uc_hierarchy} présente la hiérarchie générale des acteurs. L'acteur de base « Utilisateur » est spécialisé en quatre rôles distincts, chacun héritant de la capacité de s'authentifier. Cette hiérarchie simple reflète l'architecture RBAC du système, où chaque rôle correspond à un ensemble disjoint de permissions.

\begin{figure}[H]
    \centering
    \pumlsvg[0.6\textwidth]{uc_hierarchy}
    \caption{Hiérarchie générale des acteurs}
    \label{fig:uc_hierarchy}
\end{figure}

Le diagramme de la \ref{fig:uc_admin} détaille les cas d'utilisation de l'Administrateur. Ce rôle est le seul à pouvoir configurer les structures académiques (départements, filières, salles), gérer les comptes utilisateurs (création, modification, suppression, import CSV) et consulter les journaux d'audit. Il configure également les paramètres globaux du système (templates de jurys, durée des soutenances, configuration email).

\begin{figure}[H]
    \centering
    \pumlsvg[0.7\textwidth]{uc_admin}
    \caption{Cas d'utilisation : Administration du système}
    \label{fig:uc_admin}
\end{figure}

Le diagramme de la \ref{fig:uc_coordination} détaille les cas d'utilisation du Coordinateur, l'acteur le plus sollicité par le système. Il gère le cycle complet : des projets et groupes d'étudiants, à la constitution des jurys, en passant par la planification interactive via le Defense Designer, la validation du planning, et la génération des documents. Les relations « include » et « extend » modélisent les dépendances entre cas d'utilisation (par exemple, « Affecter un jury » inclut « Vérifier les conflits »).

\begin{figure}[H]
    \centering
    \pumlsvg[\textwidth]{uc_coordination}
    \caption{Cas d'utilisation : Coordination et Defense Designer}
    \label{fig:uc_coordination}
\end{figure}

Le diagramme de la \ref{fig:uc_jury} présente les cas d'utilisation de l'Enseignant en tant que membre de jury. Il consulte son planning personnel, déclare ses indisponibilités via un calendrier interactif, saisit les notes et observations pour les étudiants qui lui sont assignés, et accède à ses convocations. Le cas « Saisir les évaluations » inclut la consultation des informations de l'étudiant, garantissant que le jury dispose du contexte nécessaire avant de noter.

\begin{figure}[H]
    \centering
    \pumlsvg[0.7\textwidth]{uc_jury}
    \caption{Cas d'utilisation : Enseignant / Membre de jury}
    \label{fig:uc_jury}
\end{figure}

Le diagramme de la \ref{fig:uc_student} présente les cas d'utilisation de l'Étudiant. Son périmètre est le plus restreint : il consulte les informations relatives à sa soutenance (date, salle, jury), télécharge sa convocation, et gère son groupe de projet (création, recherche de coéquipiers). Le cas « Télécharger la convocation » inclut la consultation des informations de soutenance, car la convocation doit contenir la date, l'heure, la salle et la composition du jury.

\begin{figure}[H]
    \centering
    \pumlsvg[0.6\textwidth]{uc_student}
    \caption{Cas d'utilisation : Étudiant}
    \label{fig:uc_student}
\end{figure}

\subsubsection{Diagramme de classes}
\label{ch3:classes}

Le diagramme de classes traduit la structure des données persistées. On y observe une hiérarchie d'utilisateurs basée sur l'héritage JOINED (\texttt{User} avec spécialisation par rôle via le discriminant \texttt{dtype}), et les relations centrales entre les entités métier.

Plusieurs choix de conception méritent d'être soulignés. La classe \texttt{JuryMember} est une entité \texttt{@Embeddable} stockée dans une table d'association \texttt{defense\_members}, plutôt qu'une entité JPA à part entière. Ce choix simplifie la gestion des jurys en évitant une table supplémentaire tout en permettant de stocker le rôle du membre (Président, Rapporteur, Examinateur) et la référence vers l'enseignant. La classe \texttt{Evaluation} stocke le \texttt{teacherId} sous forme de \texttt{Long} plutôt que d'une relation \texttt{@ManyToOne}, ce qui permet de conserver les évaluations même en cas de modification de la composition du jury. La table d'association \texttt{defense\_session\_coefficients} (\texttt{@ElementCollection}) permet de pondérer dynamiquement les notes selon le rôle du membre du jury au sein de chaque session, offrant ainsi une flexibilité maximale dans la configuration de la notation.

\begin{figure}[H]
    \centering
    \pumlsvg[\textwidth]{class_global}
    \caption{Diagramme de classes global du système}
    \label{fig:class_global}
\end{figure}

\subsubsection{Diagrammes de séquence}
\label{ch3:sequence}

Les diagrammes de séquence modélisent les interactions entre les composants du système pour des cas d'utilisation critiques. Nous en présentons trois qui illustrent les flux les plus complexes.

Le diagramme de la \ref{fig:seq_planning} modélise le flux de planification interactive via le Defense Designer. Ce processus repose sur un cycle de validation itératif où le frontend et le backend collaborent pour garantir l'absence de conflits :

\begin{enumerate}
    \item L'utilisateur fait glisser un jury depuis la barre latérale vers une cellule du calendrier.
    \item Le hook \texttt{useDefenseSchedule} construit le contexte (\texttt{ConflictContext}) à partir des données locales.
    \item Le moteur TypeScript (\texttt{validateSlotAssignment}) exécute les 8 vérifications.
    \item En cas d'erreur, un toast est affiché avec une suggestion de résolution ; le dépôt est annulé.
    \item En cas de succès, l'état local (\texttt{schedule}) est mis à jour immédiatement (\textit{optimistic update}).
    \item Lors de la sauvegarde, les données sont envoyées au serveur qui ré-exécute les vérifications via \texttt{ConflictDetectionService} avant persistance.
\end{enumerate}

\begin{figure}[H]
    \centering
    \pumlsvg[\textwidth]{seq_planning}
    \caption{Séquence de planification interactive via le Defense Designer}
    \label{fig:seq_planning}
\end{figure}

Le diagramme de la \ref{fig:seq_auth} modélise le flux d'authentification. L'utilisateur soumet ses identifiants via le formulaire de connexion, le frontend envoie une requête POST à \texttt{/api/auth/login}, le serveur valide les credentials (vérification du hash BCrypt et du statut du compte), puis retourne un token JWT contenant l'identifiant utilisateur et son rôle. Le frontend stocke le token dans le \texttt{localStorage} et redirige l'utilisateur vers la page correspondant à son rôle.

\begin{figure}[H]
    \centering
    \pumlsvg[0.8\textwidth]{seq_auth}
    \caption{Séquence d'authentification via JWT}
    \label{fig:seq_auth}
\end{figure}

Le diagramme de la \ref{fig:seq_evaluation} modélise le flux de soumission d'une évaluation. L'enseignant saisit le score et le commentaire pour un étudiant, le frontend envoie une requête POST à l'endpoint d'évaluation, le serveur valide la note (plage 0--20), enregistre l'évaluation, met à jour le statut (PENDING $\rightarrow$ SUBMITTED), et retourne la confirmation. Le moteur de notation calculera ensuite la moyenne pondérée une fois toutes les évaluations d'un jury soumises.

\begin{figure}[H]
    \centering
    \pumlsvg[0.8\textwidth]{seq_evaluation}
    \caption{Séquence de soumission d'une évaluation}
    \label{fig:seq_evaluation}
\end{figure}

\subsubsection{Diagrammes d'activité et de machine d'états}
\label{ch3:activite}

L'algorithme de détection des conflits, moteur central du système, est modélisé par le diagramme d'activité de la \ref{fig:act_conflict}. Chaque vérification est une étape indépendante, permettant d'identifier l'ensemble des violations en un seul passage. L'algorithme exécute les 8 règles de gestion en séquence et collecte toutes les violations détectées avant de retourner le résultat au frontend.

\begin{figure}[H]
    \centering
    \pumlsvg[0.7\textwidth]{act_conflict_logic}
    \caption{Logique de validation des 8 règles de gestion (RG1-RG8)}
    \label{fig:act_conflict}
\end{figure}

Le diagramme de la \ref{fig:act_evaluation} modélise le processus d'évaluation d'une soutenance, de la saisie des notes par les membres du jury jusqu'à la clôture de la session. Les notes sont saisies individuellement par chaque membre, puis le système calcule automatiquement la moyenne pondérée en appliquant les coefficients configurés pour la session. Une fois toutes les évaluations soumises, une décision est proposée (Admis, Admis avec corrections, Ajouré) et le procès-verbal peut être généré.

\begin{figure}[H]
    \centering
    \pumlsvg[0.7\textwidth]{act_evaluation}
    \caption{Processus d'évaluation et de notation d'une soutenance}
    \label{fig:act_evaluation}
\end{figure}

Le diagramme de la \ref{fig:state_soutenance} modélise le cycle de vie complet d'une soutenance via une machine à états. La soutenance traverse successivement les états : CREEE, PLANIFIEE, JURY\_AFFECTE, CONVOQUEE, EN\_COURS, EVALUEE, et CLOTUREE. L'état ANNULE reachable depuis PLANIFIEE, CONVOQUEE et EVALUEE permet de gérer les annulations à différents stades du processus. Cette modélisation garantit que chaque soutenance suit un parcours validé et que les transitions non autorisées sont bloquées par le système.

\begin{figure}[H]
    \centering
    \pumlsvg[0.7\textwidth]{state_soutenance}
    \caption{Machine à états du cycle de vie d'une soutenance}
    \label{fig:state_soutenance}
\end{figure}

\subsection{Architecture logicielle}
\label{ch3:architecture_logicielle}

Le système adopte une architecture découplée pour maximiser la maintenabilité et la testabilité. La séparation frontend/backend via une API REST permet à chaque couche d'évoluer indépendamment.

\subsubsection{Backend (API REST)}
Le backend suit un pattern en couches strict : \textbf{Contrôleur $\rightarrow$ Service $\rightarrow$ Repository $\rightarrow$ Entité}. Cette structure permet d'isoler la logique métier (notamment le \texttt{ConflictDetectionService}) de la couche d'accès aux données. Les DTOs assurent le découplage entre les API exposées et le modèle interne. Le projet est organisé en dix packages fonctionnels :

\begin{itemize}[label=\textcolor{primary}{\textbullet}]
    \item \texttt{auth/} -- authentification JWT, mot de passe oublié, activation de compte
    \item \texttt{admin/} -- gestion des entités académiques (salles, départements, filières, configuration)
    \item \texttt{coordinator/} -- cœur métier (projets, groupes, jurys, planification, conflits, notes, documents)
    \item \texttt{teacher/} -- evaluations, indisponibilités, planning personnel
    \item \texttt{student/} -- groupe, documents, convocation, statistiques
    \item \texttt{notification/} -- système de notifications et email
    \item \texttt{common/} -- configuration Spring Security, exceptions, mappers, audit
    \item \texttt{seed/} -- initialisation des données de démonstration
    \item \texttt{user/} -- hiérarchie des utilisateurs (User, Teacher, Student, Coordinator)
\end{itemize}

\subsubsection{Frontend (SPA)}
L'interface est conçue comme une application monopage (\textit{Single Page Application}) utilisant React 19 avec TypeScript 6 en mode strict. L'état asynchrone est géré par \textbf{TanStack Query v5}, qui assure la synchronisation avec l'API, la mise en cache et le rechargement automatique. Le \textit{Defense Designer} implémente un état local optimiste pour une fluidité maximale lors du drag-and-drop, grâce à la bibliothèque \textbf{@dnd-kit/core}. Chaque page est chargée via \texttt{React.lazy()} pour optimiser le temps de chargement initial, et les routes sont protégées par des gardes \texttt{ProtectedRoute} qui vérifient le rôle de l'utilisateur avant d'autoriser l'accès.

\subsubsection{Sécurité et contrôle d'accès}
Le flux de sécurité repose sur trois mécanismes complémentaires. L'authentification utilise des tokens JWT signés avec une clé HMAC-SHA, contenant l'identifiant utilisateur et son rôle. Le filtre \texttt{JwtAuthFilter} intercepte chaque requête, extrait le token du header \texttt{Authorization} ou du cookie \texttt{jwt\_token}, le valide, et configure le contexte de sécurité Spring. Le contrôle d'autorisation est implémenté via les annotations \texttt{@PreAuthorize} sur les contrôleurs, qui restreignent l'accès aux endpoints en fonction du rôle de l'utilisateur courant. Enfin, la configuration CORS autorise les origines du frontend en développement (\texttt{localhost:5173}) et en production.

\subsection{Modélisation de la base de données}
\label{ch3:mld}

Le Modèle Logique de Données (MLD) traduit les classes JPA en tables relationnelles. La base de données H2 en développement utilise le même schéma que MySQL en production, garantissant la compatibilité entre les environnements.

L'héritage des utilisateurs est implémenté via la stratégie JOINED : la table \texttt{users} contient les attributs communs (email, password, role, firstName, lastName, isActive), tandis que les tables \texttt{teachers}, \texttt{students} et \texttt{coordinators} stockent les attributs spécifiques. Cette stratégie, plus coûteuse en jointures que la stratégie SINGLE\_TABLE, a été choisie pour sa clarté sémantique et sa facilité d'évolution.

Les données flexibles sont gérées via les annotations \texttt{@ElementCollection} : les membres d'un jury (\texttt{defense\_members}), les créneaux d'indisponibilité (\texttt{unavailability\_slots}), les coefficients d'évaluation par session (\texttt{defense\_session\_coefficients}), et les rôles d'un template de jury (\texttt{template\_roles}). Cette approche, bien que moins performante que des tables séparées pour les données volumineuses, est adaptée au volume de données attendu dans un contexte universitaire.

Les entités de configuration (GeneralSettings, DefenseSettings, EmailConfig, DocumentConfig) utilisent un pattern singleton avec un identifiant fixe (\texttt{id=1L}), simplifiant la gestion des paramètres globaux par le système.

\begin{figure}[H]
    \centering
    \pumlsvg[\textwidth]{mld_schema}
    \caption{Modèle Logique de Données (MLD) synchronisé}
    \label{fig:mld_schema}
\end{figure}

\subsection{Normes et standards respectés}
\label{ch3:normes}

Le développement du système suit des normes et standards industry pour garantir la qualité, la maintenabilité et la cohérence du code source.

Côté Java, le code respecte les conventions de nommage Java (camelCase pour les méthodes et variables, PascalCase pour les classes), les directives de codage Google (limitation à 100 caractères par ligne, gestion des imports non utilisés), et les bonnes pratiques Spring (injection de dépendances par constructeur, annotations de validation sur les DTOs). Les noms de tables et de colonnes suivent la convention snake\_case en base de données, reflétée par les annotations \texttt{@Table} et \texttt{@Column}.

Côté TypeScript, le code suit les conventions Airbnb (formatage Prettier, règles ESLint strictes), le typage strict est activé (\texttt{strict: true} dans le \texttt{tsconfig.json}), et les noms de fichiers utilisent le kebab-case (\texttt{conflict-engine.ts}, \texttt{use-defense-schedule.ts}). Les composants React suivent le pattern fonctionnel avec hooks, sans/classes.

Les messages de commit suivent les conventions Git : préfixes sémantiques (\texttt{feat:}, \texttt{fix:}, \texttt{chore:}, \texttt{docs:}), description en anglais, imperative mood. Les issues GitHub documentent les fonctionnalités et les corrections, avec des labels de priorité et de catégorie.

\subsection{Conclusion}
\label{ch3:conclusion}

Ce chapitre a permis de formaliser les besoins du système et de concevoir une architecture capable de supporter la complexité de la planification académique. La définition rigoureuse des 8 règles de gestion, la modélisation UML complète (cas d'utilisation, classes, séquence, activité, machine à états) et la description de l'architecture logicielle en couches servent désormais de base technique pour l'implémentation détaillée qui sera présentée au chapitre suivant. Le respect des normes de codage et des standards industry garantit la maintenabilité du système à long terme.
